\documentclass[../../main2.tex]{subfiles}

\begin{document}

	% ======================================================================
	% 骨架文件：小节标题与追问链为终版；正文由后续会话填充。
	% 扩增模式标注：本章 = 有压扩增（稀缺约束）。
	% 原书材料接入点：
	%   - 原书 part3 ch1（经济基础）→ 本章主体，改由「文化+稀缺」两步推导导出
	%   - §10.1 是本 Part 的枢纽节（notes/00-architecture.md 明示「必须完整写」）
	% ======================================================================

	\chapter{通道二：经济（模因 + 稀缺约束）}
	\label{ch:economy}

	\begin{motivation}
		上一章遗留的问题：模因传播不消耗原体、也就无价——但现实里大量传播必须付费。如果一种模因的传播\textbf{必须}消耗稀缺资源，它还是模因吗？\\
		\textbf{核心动机}：用「加入稀缺约束」这一步，从文化模因中推导出经济——供求不是经济学的基本假设，而是「模因 + 稀缺」两步推导的产物。\\
		\textbf{扩增模式}：有压——复制变转移，总量守恒，零和开始起作用。
	\end{motivation}

	\begin{logicmap}
	\begin{tikzpicture}[node distance=0.9cm, every node/.style={draw, rectangle, rounded corners, align=center}]
	\node (q) {传播要花钱，还是模因吗？};
	\node (s) [below=of q] {加入稀缺约束（§10.1，枢纽节）};
	\node (sd) [below=of s] {稀缺→供方；偏好→需求（§10.2）};
	\node (h) [below=of sd] {追涨杀跌的理性化（§10.3）};
	\node (c) [below=of h] {经济的贡献分解（§10.4）};
	\node (n) [below=of c] {交换真的等价吗？（抛给第11章）};
	\draw[-{Stealth}] (q) -- (s);
	\draw[-{Stealth}] (s) -- (sd);
	\draw[-{Stealth}] (sd) -- (h);
	\draw[-{Stealth}] (h) -- (c);
	\draw[-{Stealth}] (c) -- (n);
	\end{tikzpicture}
	\end{logicmap}

	% ==================================================================
	\section{加入稀缺约束}
	\label{sec:econ-scarcity}

	\textcolor{gray}{\textit{【骨架 · 追问链（终版）】复制变成转移 → 无压扩增变有压扩增 → 总量守恒。这是本 Part 的枢纽节，必须完整写。守恒带来了什么新性质？→ 零和博弈开始起作用：你多得我就少得。}}

	\textcolor{gray}{（正文待写：标准句式完整展开——「如果传播不消耗原体（你把想法告诉我，你还有这个想法），扩张是无限的。但现在加入约束：传播需要消耗稀缺资源。这一步改变了所有性质——这就是经济行为从文化行为中分离出来的时刻。」复制 $A \to A+B$ vs 转移 $A \to B$；价值不凭空创造或消灭。）}

	\begin{readingnote}
		\textcolor{gray}{（待写：你有没有想过，为什么「分享」这个词在朋友圈免费，在股市要收印花税？）}
	\end{readingnote}

	\paragraph*{逻辑衔接}
	\textcolor{gray}{（待写）守恒性落地之后，最古老的经济学图景自己长出来了。下一节「稀缺→供方；偏好→需求」做两步推导。}

	% ==================================================================
	\section{稀缺到供方，偏好到需求}
	\label{sec:econ-supplydemand}

	\textcolor{gray}{\textit{【骨架 · 追问链（终版）】稀缺产生稀缺端（供方），需求来自模因层面的偏好（第2章的目标函数在此上岗）。价格从哪来？→ 稀缺信号——价格是压缩进一个数字的供需信息，且会失真。}}

	\textcolor{gray}{（正文待写：供求不是假设而是推导；价格 = 信息压缩器（呼应原书变长编码）；外部性与市场失灵是压缩失真的形式。）}

	\begin{questionbox}
		\textcolor{gray}{（待写：反驳位——「供求定律有反例（吉芬商品、炫耀性消费），推导还有效吗？」回答：反例不推翻「两步推导」，只说明偏好维度更复杂——恰好回到第2章的多目标。）}
	\end{questionbox}

	\paragraph*{逻辑衔接}
	\textcolor{gray}{（待写）价格是信号，但信号有噪声。噪声下的行为是理性的还是疯狂的？下一节「追涨杀跌的理性化」沿用原书答案。}

	% ==================================================================
	\section{追涨杀跌的理性化}
	\label{sec:econ-herding}

	\textcolor{gray}{\textit{【骨架 · 追问链（终版）】沿用原书：追涨杀跌是不完全信息下的启发式——模仿他人作为信号处理。这是错误还是理性？→ 是理性：当私有小、公共信息在价格里时，跟风是信息约束下的最优近似。}}

	\textcolor{gray}{（正文待写：信息级联（information cascade）的最小模型；泡沫与崩盘作为级联的两端；「理性个体的理性加总出集体非理性」。）}

	\begin{readingnote}
		\textcolor{gray}{（待写）}
	\end{readingnote}

	\begin{reader_anticipation}
		\item \textcolor{gray}{（待写）级联什么时候断裂？→ 第11章 §11.2：认同的凭空增减——政治守恒性弱于经济的根源。}
	\end{reader_anticipation}

	\paragraph*{逻辑衔接}
	\textcolor{gray}{（待写）经济通道的机制齐了。它在结果里占多少？下一节「经济的贡献分解」——并直视一个常见的夸大。}

	% ==================================================================
	\section{经济的贡献分解}
	\label{sec:econ-contribution}

	\textcolor{gray}{\textit{【骨架 · 追问链（终版）】$C_{\text{经济}}$ 的代入计算。诚实说明：经济贡献常被高估，因为价值是「测量的」不是「物理的」——价格会失真，GDP 不是自然常数。怎么测？→ 用价格，但必须并列价格失真的方向与幅度。}}

	\textcolor{gray}{（正文待写：贡献表示例；被高估与被低估的典型情形（家务劳动、生态成本）；\lowfreq 标注测量边界。）}

	\paragraph*{逻辑衔接}
	\textcolor{gray}{（待写）经济交换在课本里是等价交换——双方地位对称。真的吗？下一章的入口就在这里：一旦交换双方的身份变得重要，等价交换就变形了。}

	% ==================================================================
	\section{本章小结与衔接}
	\label{sec:econ-summary}

	\textcolor{gray}{（正文待写：收束——经济 = 加了稀缺约束的模因行为；守恒、零和、价格信号、级联。嵌套预告：文化 $\supset$ 经济。）}

	\paragraph*{逻辑衔接}
	\textcolor{gray}{（待写）第11章「通道三：政治（经济 + 身份权重）」：当交换中加入不可货币化的偏好——人情、义务、权力——政治从经济中分离出来。}

\end{document}
